\documentclass[11pt]{article}
\usepackage[margin=1in]{geometry}
\usepackage{parskip}
\usepackage{xcolor}
\usepackage{titlesec}
\titleformat{\section}[block]{\large\bfseries\color{blue!50!black}}{}{0pt}{}
\setlength{\parskip}{0.5em}

\title{Speaker script (final version) for\\
       \emph{Deep Learning for Mean Field Games with Non-Separable Hamiltonians}}
\author{Companion to \texttt{presentation\_mfdgm.tex}\\
        Hikmatullo Ismatov}
\date{}

\begin{document}
\maketitle

\textit{Each numbered section below corresponds to one slide of the final
presentation file.  Read each section as one ``unit of speech,'' roughly
60--90 seconds per slide.  Total runtime at a calm pace: about 35--40 minutes
of content plus questions.  Section pages produced automatically by Metropolis
between parts are not counted here; on those, just say ``Part X'' and the
name of the section in a single sentence.}

% ================================================================
\section*{Slide 1 --- Title page}

Hello everyone, thank you for being here today.

I am Hikmatullo Ismatov, a graduate student at KAUST, working with
Professor Diogo Gomes.

Today I will walk us through a paper by Mouhcine Assouline and Badr
Missaoui, published in 2023 in \emph{Chaos, Solitons and Fractals}.  It is
called \emph{Deep Learning for Mean Field Games with Non-Separable
Hamiltonians}.

I will explain everything slowly, define every technical term, and at the
end I will connect this paper to a new method we are developing in our
group.

% ================================================================
\section*{Slide 2 --- What this talk covers}

The talk has nine short parts.  We begin with foundations: what is a Mean
Field Game.  Then we ask why these games are hard to solve in high
dimensions.

Once we understand the problem, we will see that deep solvers come in two
big families: GAN-based methods and physics-informed methods.  Before we
read the paper, we need to understand three background tools: what a GAN
is, what the Hopf formula is, and what APAC-Net is.  These three sit
behind everything we will discuss.

Then we read the paper itself, look at its convergence theory and its
experiments.  And at the very end I will briefly tell you about our
follow-up work, which builds on this paper.

Please interrupt at any time with questions.

% ================================================================
\section*{Slide 3 --- What is a Mean Field Game (intuition)}

Imagine a thousand drivers on the same highway.  Each driver wants to
reach their destination quickly.  But on crowded sections, everyone slows
down; on empty sections, everyone speeds up.

So each driver's best strategy depends on what the rest of the population
is doing.  And the rest of the population is doing whatever they think is
best for them.  This is circular.

A Mean Field Game is the mathematical theory that resolves this
circularity.  It was developed independently by Lasry and Lions in 2007,
and by Huang, Caines, and Malham\'e in 2006.  The idea is to find a Nash
equilibrium: a strategy from which no single agent benefits by changing
their mind alone.

The phrase \emph{mean field} comes from physics.  It means that no
individual agent is large enough to matter on its own; only the average
behaviour of the whole crowd matters.

% ================================================================
\section*{Slide 4 --- The MFG system}

At Nash equilibrium, the game is described by two coupled partial
differential equations.

The first one is the Hamilton-Jacobi-Bellman equation, or HJB.  It governs
the value function $\varphi$, and it runs backward in time, from the
terminal time $T$ back to time zero.  Each agent uses $\varphi$ as their
internal optimal cost.

The second equation is the Fokker-Planck equation, or FP.  It governs the
population density $\rho$, and it runs forward in time, from the initial
density at time zero.

The boundary conditions are simple.  The initial density is given, and
the terminal cost is also given, possibly depending on the terminal
density.

Look at the small diagram at the bottom.  The two equations are coupled:
the gradient of $\varphi$ tells agents where to go, which becomes the
drift in the Fokker-Planck equation.  And the density $\rho$ enters the
cost in the HJB.  Neither equation can be solved alone.

The parameter $\nu$ is the diffusion: $\nu = 0$ is the deterministic
case, and $\nu > 0$ adds random noise.

% ================================================================
\section*{Slide 5 --- Separable vs.\ non-separable Hamiltonians}

This is the single most important distinction in the talk.  Please pay
close attention.

A Hamiltonian is called \emph{separable} if it splits as $H_0$ of $x$, $p$
minus $f_0$ of $x$, $\rho$.  The momentum part and the density part are
in different terms and never multiplied together.  A typical example is
one half $p$ squared minus lambda times $\rho$.

A Hamiltonian is called \emph{non-separable} if $\rho$ and $p$ appear
mixed, that is multiplied together.  The canonical example, on the right
of the slide, is the traffic flow model of Huang and collaborators:
one half $p$ squared minus the term one minus $\rho$ times $p$.  Density
and momentum are entangled here.

Why does this matter?  Because some deep-learning methods, specifically
the GAN family, only work when $H$ is separable.  When $\rho$ and $p$ are
entangled, those methods break.  The traffic flow model is the canonical
hard example that the paper of Assouline and Missaoui addresses.

% ================================================================
\section*{Slide 6 --- The curse of dimensionality}

So why don't we just solve the MFG with classical numerical methods?
Because they die in high dimensions.

Classical methods, like finite differences, put down a grid of points
covering the domain.  Look at the table on the slide.  In one dimension,
one hundred grid points.  In two dimensions, ten thousand.  In three
dimensions, one million.  Doable.

In ten dimensions: ten to the twentieth.  Already hard.

In one hundred dimensions: ten to the two hundredth.  This number has
two hundred digits.  It is more than the number of atoms in the entire
universe.

This is called the curse of dimensionality.  Classical grid-based methods
simply cannot handle high-dimensional MFGs.

But many MFG applications, like multi-agent reinforcement learning and
financial markets, live in hundreds of dimensions.  We need a different
approach.

% ================================================================
\section*{Slide 7 --- Neural networks to the rescue}

A neural network is just a parameterised function.  The parameters are
called weights, and we can tune them by gradient descent.

There is a classical theorem by Hornik from 1991, the Universal
Approximation Theorem.  It says that a neural network with a single
hidden layer and enough neurons can approximate any continuous function
to any desired accuracy, and this works in any dimension, without a grid.

So the strategy is: replace the unknown functions $\varphi$ and $\rho$ by
neural networks, then optimise the weights so the networks solve the PDE.

The key point: the number of parameters of a neural network does not
explode exponentially with dimension.  A hundred-dimensional network with
a hundred neurons per layer has only about ten thousand parameters.
That is tractable on a laptop.  This is why deep learning is the right
tool for high-dimensional MFGs.

% ================================================================
\section*{Slide 8 --- Two ways to train the networks}

There are two big families of methods.

On the left of the slide, the physics-informed family.  Sometimes called
PINN, sometimes DGM.  The idea is to plug the network directly into the
PDE and minimise the residual.  At random collocation points, evaluate
the network and its derivatives, plug into the PDE, square the error,
average, gradient-descend.  This works for any Hamiltonian.  Today's
paper sits in this family.

On the right, the variational or GAN-based family.  The idea is to
rewrite the PDE as a saddle-point problem and let two networks play a
game.  Examples are APAC-Net by Lin and collaborators in 2021, and
MFGAN by Ruthotto and collaborators in 2020.  These methods only work
when the Hamiltonian is separable, because they need the Hopf formula to
turn the PDE into a saddle-point problem.

At the end of the talk I will tell you about our work, which bridges
these two columns.

% ================================================================
\section*{Slide 9 --- Generative models: the goal}

To understand what a GAN is, forget about MFGs for a moment.

Here is the goal of a generative model.  I show you a million photos of
cats.  I ask you to write a computer program that, when I press a button,
outputs a new photo of a cat that I have never seen, but it has to look
real.

The mathematical version: you are given samples from an unknown
probability distribution $P_{\text{data}}$.  You want to build a program
that produces fresh samples from that distribution.

This is the generative modelling problem.  Generative Adversarial
Networks, or GANs, are one solution.

% ================================================================
\section*{Slide 10 --- GAN: a two-player game}

A GAN, introduced by Goodfellow and collaborators in 2014, consists of
two neural networks playing a game.

The first network is the \emph{generator}, denoted $G_\theta$.  It takes
random noise $z$ as input and produces a fake sample $x$.  Its goal is
to make the fakes look indistinguishable from the real samples.

The second network is the \emph{discriminator}, denoted $D_\omega$.  It
takes a sample, real or fake, and outputs the probability that the
sample is real.  Its goal is to correctly catch the fakes.

The diagram on the right of the slide shows the flow: noise enters the
generator, produces a fake; the discriminator sees both real and fake,
and tries to say which is which.

The two networks compete.  The objective at the bottom of the slide says:
the discriminator wants to maximise the probability of correct
classification.  The generator can only influence the second term, and
tries to minimise it, that is, to fool the discriminator.

At the Nash equilibrium of this game, if both networks are powerful
enough, the generator produces samples exactly from the true data
distribution.

Notice: this is itself a Nash equilibrium of a two-player game.  That is
the first hint that GANs and MFGs are linked.

% ================================================================
\section*{Slide 11 --- The Wasserstein GAN}

The original GAN has practical problems.  Training can be unstable.
There is mode collapse, where the generator produces only one or two
types of samples.  There are vanishing gradients.

So in 2017, Arjovsky, Chintala, and Bottou proposed a better version:
the Wasserstein GAN.

In a Wasserstein GAN, instead of asking the discriminator real or fake,
we ask it to assign a score to every sample.  The discriminator is
restricted to be one-Lipschitz, meaning it cannot change too fast.  The
generator tries to make its samples score high.

There is a beautiful theorem from optimal transport, called the
Kantorovich-Rubinstein duality.  It says: the maximum value of this game
equals the Wasserstein-one distance between the real and generated
distributions.

This means training a Wasserstein GAN is literally minimising the
Wasserstein distance between fake and real.  The training loss has a
clean geometric meaning at every step.  That is why W-GANs are easier
to train and behave better.

% ================================================================
\section*{Slide 12 --- Why are GANs and MFGs structurally linked?}

Now look at the two formulas on the slide.

The first one is the Wasserstein GAN: a min-max problem.  The generator
minimises, the discriminator maximises.

The second one is the saddle-point form of a separable MFG, derived by
Cao, Guo, and Lauri\`ere in 2021.  Again a min-max problem.  Here the
density $\rho$ plays the role of the generator, and the value function
$\varphi$ plays the role of the discriminator.

The two have exactly the same structure: a distribution-producing object
plays against a function being optimised.  This shared structure is the
reason GAN-based methods can solve separable MFGs.

But, very important: deriving this saddle-point form requires the Hopf
formula, and the Hopf formula requires separability.  We will see why on
the next two slides.

% ================================================================
\section*{Slide 13 --- The Hopf formula}

The Hopf formula is a way to write the solution of a Hamilton-Jacobi
equation as an optimisation problem.

For the simplest HJ equation, $\partial_t \varphi$ plus $H$ of gradient
$\varphi$ equals zero, the Hopf formula on the slide says: $\varphi$ at
time $t$ and position $x$ equals the sup over $y$ of the inf over $z$ of
an explicit expression involving the initial data, the time, and the
Hamiltonian.

The right-hand side is a min-max problem.

This is exactly the structure GAN training can attack.  In one sentence:
the Hopf formula is the bridge that turns an MFG into a GAN.

% ================================================================
\section*{Slide 14 --- Why the Hopf formula needs separability}

For an MFG, the Hamilton-Jacobi equation involves $H$ of $x$, $\rho$,
gradient $\varphi$.

If $H$ is separable, then $f_0$ of $x$, $\rho$ is constant in $p$.  You
can pull it out and treat it as a known source term.  The Hopf formula
then applies to the remaining $H_0$ of $x$, $\nabla\varphi$ piece.  You
get a clean min-max formulation.  This is exactly what APAC-Net uses.

If $H$ is non-separable, density and momentum are entangled.  The
Legendre transform variable cannot factor cleanly.  The Hopf formula
does not give a min-max reformulation any more.

This is the structural barrier.  The whole reason APAC-Net and MFGAN
cannot solve the traffic flow MFG.

Notice this is not a numerical problem.  It is an algebraic one.

% ================================================================
\section*{Slide 15 --- APAC-Net}

APAC-Net was introduced by Lin and collaborators in 2021.  It targets
the separable case.

APAC-Net has two neural networks.  The first is a generator $G_\theta$,
which takes random noise $z$ and time $t$, and outputs a position $x$
in $d$-dimensional space.  As you sample $z$ many times at a fixed $t$,
you get many positions, and these positions distribute according to the
population density at that time.  So the generator parameterises the
density distributionally.

The second network is a critic $\varphi_\omega$.  It takes time and
position and outputs a scalar.  This is the value function.

The objective comes from the Hopf reformulation: a min-max where the
generator minimises and the critic maximises.

Training is exactly the same as a Wasserstein GAN: alternating Adam
steps on the two networks.

% ================================================================
\section*{Slide 16 --- What APAC-Net gives and what it does not}

What APAC-Net gives you, on the positive side:

You get sampling access to the density through the generator.  This is
the distributional representation we keep mentioning.  You get the
value function directly.  You have a Wasserstein-one convergence
guarantee, by W-GAN theory.  And the method scales to high dimensions.

What APAC-Net cannot do:

It cannot handle non-separable Hamiltonians.  Period.  In particular,
the traffic flow MFG-LWR is out of reach.  Also, no explicit convergence
rate has been proved.

The small comparison table on the slide summarises this.  Now the
natural question: can we handle non-separable Hamiltonians with a deep
solver?  That is what MFDGM addresses, on the next slide.

% ================================================================
\section*{Slide 17 --- MFDGM: the idea}

Here is the key insight of Assouline and Missaoui.

Stop trying to reformulate the MFG as a min-max problem.  Just plug the
networks into the PDEs and minimise the residual directly.

There are two networks, same as before: a value network $\varphi_\omega$
and a density network $\rho_\theta$.

If they truly solve the MFG, then by definition the PDE residuals are
zero.  So define loss functions that measure how far the residual is
from zero, and let an optimiser fix the weights.

The crucial point: the Hamiltonian $H$ appears \emph{directly inside}
the loss.  We never reformulate as a saddle-point.  No Hopf formula is
needed.  So this method works for any Hamiltonian, separable or not.

% ================================================================
\section*{Slide 18 --- The loss functions of MFDGM}

There are two loss functions, one per equation.

The HJB loss measures how badly the value network violates the HJB
equation.  At each random collocation point in the domain, we plug the
two networks into the HJB and square the residual.  We also add a
penalty for not satisfying the terminal cost.

The FP loss does the same for the Fokker-Planck equation: random
points, plug in, square the residual, plus an initial-condition penalty.

The points $(t_b, x_b)$ are sampled randomly inside the domain.

This is exactly the PINN paradigm of Raissi, Perdikaris, and
Karniadakis from 2019.  Differentiate the network using automatic
differentiation, plug into the PDE, average the residual over random
collocation points, gradient-descend.

% ================================================================
\section*{Slide 19 --- Algorithm 1 of MFDGM}

The training algorithm is exactly what you would expect from the loss
functions.

Initialise both networks with random weights.  Repeat $K$ times: sample
random points, compute the HJB loss using the current networks, take
one Adam step on the value network.  Then sample fresh random points,
compute the FP loss, take one Adam step on the density network.

This is alternating training.  The HJB step uses the current density;
the FP step uses the just-updated value function.

The authors emphasise that they use two separate Adam optimisers, one
per network, rather than a single combined loss.  In their experiments
this converges faster than the earlier DGM-MFG method, which used a
single combined loss.

% ================================================================
\section*{Slide 20 --- Theorem 3.1 (existence of approximating networks)}

Now the convergence theory.

Theorem 3.1 of the paper is an existence statement.  Under regularity
conditions H1, H2, H3, for any positive epsilon, there exist neural
networks whose residual losses are at most a constant times epsilon.

What does this mean?  Networks exist that make the residual arbitrarily
small.

This is a statement about \emph{expressivity}, not training.  It says
the network class is rich enough.  It does not say we can find these
networks by gradient descent.  The proof uses the Universal
Approximation Theorem in a version that controls derivatives, not just
function values.

% ================================================================
\section*{Slide 21 --- Theorem 3.2 (convergence of trained networks)}

Theorem 3.2 is the main convergence result.  Under stronger conditions
H4, H5, H6, the trained networks converge to the true MFG solution, in
$L^p$ for any $p$ less than two.

Stop and unpack this.

The trained networks converge.  This is the good news.

The convergence is in $L^p$ with $p$ strictly less than two.  This is a
weak statement.  It is weaker than $L^2$, much weaker than $L^\infty$,
and weaker than Wasserstein.  It permits transient spikes in the
density, in particular.

The theorem does not give an explicit rate.  Just convergence as the
network width goes to infinity.

So this is a positive result, but a soft one.  It is exactly the gap
our follow-up work addresses.

% ================================================================
\section*{Slide 22 --- Test 1: analytic separable benchmark}

Now the experiments.  First, a separable test case with a known
analytic solution, taken from the APAC-Net paper.

The Hamiltonian and the terminal cost are chosen so that the true
solution is explicit: $\varphi$ equals one half $x$ squared minus
$d$ times $t$, and $\rho$ is a standard Gaussian.

The paper trains MFDGM in one dimension, three hidden layers of one
hundred neurons each, using ResNet architecture, ten thousand
iterations.

The result, in Figure 2 of the paper, shows the MFDGM solution overlaid
on the analytic solution at three time slices.  The two curves are
essentially indistinguishable.  The $L^2$ relative error drops to about
one percent.

This is a sanity check, and it passes.

% ================================================================
\section*{Slide 23 --- Comparison with APAC-Net, MFGAN, DGM-MFG}

The paper also compares MFDGM against three other deep solvers on the
same separable benchmark.

Figure 9 of the paper shows the relative error of the density and the
value function for all four methods over iterations.

MFDGM achieves the lowest error.  MFGAN is comparable.  APAC-Net is
slightly limited because its density is represented distributionally,
and to compare to a pointwise solution they had to estimate it via
kernel density estimation.  DGM-MFG is slower because it uses a single
combined loss instead of two separate losses.

A caveat: this comparison is on a separable benchmark, so APAC-Net and
MFGAN are allowed to play.  The real test is the next slide.

% ================================================================
\section*{Slide 24 --- The big test: MFG-LWR traffic flow}

This is the experiment that justifies the entire paper.

The problem is the MFG-LWR traffic flow with non-separable Hamiltonian
$H$ equals one half $p$ squared minus $1-\rho$ times $p$.  Domain is
one-dimensional, time horizon one, terminal cost zero, and the initial
density is a Gaussian bump.

APAC-Net cannot run on this problem.  The Hopf formula breaks.  Only
MFDGM and classical Newton methods can.

The paper runs MFDGM both deterministically with $\nu$ equals zero and
stochastically with $\nu$ equals zero point five.  The figure on the
right shows density, optimal cost, and speed at different times.  The
solution evolves consistently with traffic-flow intuition: the density
bump disperses, the speed adjusts.

This is the unique achievement of MFDGM: a deep solver for non-separable
MFGs that scales to high dimensions.

% ================================================================
\section*{Slide 25 --- Where MFDGM leaves a gap}

Now let me describe what our group is doing on top of this paper.

Look at the comparison table on the slide.  Columns: handling of
non-separable Hamiltonians, representation of the density, convergence
metric, and whether there is an explicit rate.

APAC-Net handles only separable, but has distributional $\rho$ and
Wasserstein-one convergence.  No explicit rate.

MFDGM handles any Hamiltonian, the good news.  But its $\rho$ is a
pointwise neural network.  Its convergence is $L^p$ for $p$ less than
two, which is weak.  And no explicit rate.

So neither method gives you all three properties: distributional $\rho$,
Wasserstein convergence, and an explicit rate.

Our question is: can we close all three gaps simultaneously, at least
in some useful regime, without giving up non-separable handling?

% ================================================================
\section*{Slide 26 --- Our insight: freeze the density}

Our key insight is very simple.  Here it is.

If $\rho$ is held fixed, any non-separable Hamiltonian becomes
separable.

Look at the diagram.  The full non-separable MFG is on the left.  We
freeze the density at the current iterate $\rho^k$.  Now the effective
Hamiltonian $H^k$ of $x$ and $p$ depends only on $x$ and $p$, because
$\rho^k$ is just a fixed function.  It is separable in the sense the
GAN solvers require.

So we apply APAC-Net to this separable sub-problem.  APAC-Net returns a
new density $\rho^{k+1}$.  We update, and repeat.

This is essentially a fixed-point iteration on the density.  Each step
involves one APAC-Net call as an inner solver.  We call the whole
scheme Extended-MFGAN.

% ================================================================
\section*{Slide 27 --- Weakly non-separable Hamiltonians}

To analyse this rigorously, we need to define the right class of
Hamiltonians.

We say $H$ is \emph{weakly non-separable with strength $\eps$} if it
decomposes into three parts.  First, a strongly convex base $H_0$ with
convexity constant $c_0$.  Second, a Lasry-Lions monotone interaction
$f_0$ with constant lambda.  Third, a coupling residual $R$, scaled by
$\eps$, that is Lipschitz in both $\rho$ and $p$.

Special cases.  When $\eps$ is zero, you recover the standard separable
MFG.  When $R$ equals $\rho$ times $p$, you recover the traffic flow
model.  When $R$ depends only on $\rho$, the momentum-dependent constants
vanish.

The most important quantity is the boxed condition at the bottom:
$\eps$ times $L_R$ strictly less than lambda.  We call this the
\emph{GAN-tractable regime}.  Below the threshold, our method
provably contracts.  Above the threshold, it can fail.

The threshold balances the coupling strength against the restoring
force from the monotone interaction.

% ================================================================
\section*{Slide 28 --- The Extended-MFGAN algorithm}

The algorithm is the obvious thing.

Initialise the density.  Then loop.  At each step, build the effective
Hamiltonian by plugging in the current density.  Run an inner GAN
solver, for example APAC-Net, on the separable MFG.  Update the
density, repeat.

A nice sanity check: when $\eps$ is zero, the effective Hamiltonian
does not depend on the current density.  The outer loop converges in
one step, and we recover APAC-Net exactly.

The total cost is $K$ outer iterations times one APAC-Net call.  Once
we have proved the contraction rate $\kappa$, the number of outer
iterations needed to reach accuracy delta is explicit: log of one over
delta, divided by log of one over $\kappa$.

% ================================================================
\section*{Slide 29 --- Comparison: where we sit}

This is the punchline slide.

Compare the three methods.

APAC-Net: separable only, distributional $\rho$, Wasserstein-one
convergence, no rate.

MFDGM: handles non-separable, but pointwise $\rho$, $L^p$ for $p$ less
than two, no rate.

Our Extended-MFGAN: handles non-separable in the weakly non-separable
regime, distributional $\rho$ through the inner APAC-Net, $L^2$ and
Wasserstein-two convergence, with explicit rate $\kappa^k$.

We close all three gaps simultaneously, within the weakly non-separable
regime.

The trade-off is honest.  We are restricted to the condition
$\eps L_R$ strictly less than lambda.  Above that threshold, MFDGM
remains the right tool.  The two methods are complementary, not
competing.

% ================================================================
\section*{Slide 30 --- Thank you}

This brings me to the end of the talk.

In one sentence: today we walked through a paper that pushes deep
learning for Mean Field Games into the non-separable regime, and I
showed how our follow-up work strengthens its guarantees to
Wasserstein-two convergence with an explicit and sharp convergence
rate.

Thank you for your attention.  I am happy to take any questions.

\end{document}
